% ============================================================
%  Chapter 8 — Apps: The Workers Inside Your Phone
%  Inside a Smartphone: The Tiny World in Your Hand
% ============================================================

\chapter{Apps: The Workers Inside Your Phone}

\chapteropening{%
  The hardware — the processor, screen, and sensors — is like an empty factory.
  Apps are the workers who give that factory a purpose.
}
\chapterbadge{Mission: Put the Workers to Work!}

% --- Opening ---
\section*{Meet Appy}

\character{Appy}{Hi! I am Appy, and I keep busy every second of every day.
  I am a programme — a set of instructions that tells Chip exactly what to do.
  Without me, your phone's hardware would just sit there doing nothing.
  Every game, map, music player, and camera programme — that is an app like me!}

\sectionrule

% --- What apps do ---
\section*{What Apps Do}

An \term{app} (short for \emph{application}) is a programme designed to do
a specific job on a smartphone.
Apps are written by developers — people who write instructions in code.

Apps use the phone's hardware to do their work:
\begin{itemize}
  \item A maps app uses the GPS sensor, screen, and internet connection.
  \item A camera app uses the camera sensor, touchscreen, and storage.
  \item A music app uses the storage, processor, and speaker.
\end{itemize}

\begin{picturethis}
  Think of the phone's hardware as a well-equipped kitchen:
  it has an oven, fridge, knives, and pans.
  An app is the chef who decides what to cook and how to use every tool.
  Without the chef, the kitchen does nothing.
  Without the kitchen, the chef has nothing to work with.
\end{picturethis}

\sectionrule

% --- The operating system ---
\section*{The Operating System: The Manager}

Apps do not work alone. They all depend on the \term{operating system} (OS).

The operating system (like Android or iOS) is the main programme that:
\begin{itemize}
  \item Manages all the hardware (screen, processor, battery).
  \item Lets multiple apps run without getting in each other's way.
  \item Controls which apps are allowed to use which sensors or data.
  \item Provides the home screen, notifications, and settings.
\end{itemize}

Apps are built on top of the operating system — they cannot bypass it.
That keeps the phone organised and helps protect your data.

\sectionrule

% --- Types of apps ---
\section*{Apps for Every Purpose}

\begin{quickmap}
  \textbf{Apps by category:}
  \begin{center}
  \begin{tabular}{ll}
    \textbf{Category} & \textbf{Examples} \\
    \hline
    Communication   & Messaging, email, video calls \\
    Entertainment   & Games, music, video streaming \\
    Learning        & Educational, language, maths \\
    Navigation      & Maps, GPS directions \\
    Creativity      & Drawing, photo editing, music making \\
    Productivity    & Notes, calendars, to-do lists \\
    Health          & Step counter, sleep tracker \\
  \end{tabular}
  \end{center}
\end{quickmap}

\sectionrule

% --- Installing and updating ---
\section*{Installing and Updating Apps}

Apps are downloaded from an \term{app store} — a digital shop where developers
post their apps for people to install.
Some apps are free to download but make money through adverts or extra paid features.
That is why it is smart to ask a trusted adult before buying anything in an app.

A \term{update} is a new version of an app that:
\begin{itemize}
  \item Fixes bugs (small errors in the code).
  \item Adds new features.
  \item Improves security.
\end{itemize}

Keeping apps updated is important — old versions can have security weaknesses.

\begin{funfact}
  As of 2024, there are more than \textbf{3.5 million apps} available on
  the Google Play Store alone. Developers around the world upload
  thousands of new apps every single day.
\end{funfact}

\sectionrule

\begin{newwords}
  \begin{description}[labelwidth=4cm, leftmargin=4.2cm]
    \item[\term{App}]               A programme that runs on a phone to do a specific job.
    \item[\term{Operating system}]  The main programme managing all hardware and apps (iOS, Android).
    \item[\term{Update}]            A new version of an app or OS that fixes bugs and adds features.
    \item[\term{App store}]         A digital shop where apps are downloaded and installed.
  \end{description}
\end{newwords}

\sectionrule

\begin{tryit}
  \textbf{App audit!}

  With a grown-up's help, look at the apps installed on a smartphone.
  \begin{enumerate}[label=\textbf{\arabic*.}]
    \item Count how many apps are installed.
    \item Group them into categories (games, learning, communication, etc.).
    \item Which category has the most apps?
    \item Are there any apps that have never been opened? Should they be deleted?
  \end{enumerate}
\end{tryit}

\sectionrule

\begin{bigidea}
  An \textbf{app} is a programme that tells the phone's hardware what to do.
  The \textbf{operating system} manages all apps and hardware.
  Keeping apps \textbf{updated} keeps the phone working well and safely.
\end{bigidea}

\character{Appy}{I can help you find your way, but only because of one
  very special helper high above the Earth. Meet Orbi in Chapter 9!}
